%!TEX root = main.tex
\section{Introducción}
\PARstart{E}{l} procesamiento de imágenes digitales en tiempo real es una piedra angular de aplicaciones modernas que van desde la visión artificial en vehículos autónomos hasta la mejora de imágenes médicas. Una de las operaciones más fundamentales y computacionalmente intensivas es el suavizado de imágenes, comúnmente utilizado para la reducción de ruido mediante filtros espaciales como el filtro de caja (box filter). A medida que la resolución de las imágenes aumenta, el procesamiento secuencial tradicional se vuelve insuficiente, lo que motiva la implementación de paradigmas de paralelismo multinúcleo y heterogéneo.

Este trabajo explora la eficacia de OpenMP no solo para la explotación de hilos en la CPU, sino también para el "offloading" de cómputo hacia aceleradores gráficos (GPU). El objetivo es documentar las ventajas y desventajas de cada aproximación, analizando cómo factores como la gestión de memoria y el overhead de inicialización impactan en el rendimiento real de un kernel de suavizado aplicado a una imagen fractal de alta resolución.